\section{Composition of Linear Maps}
\bx{
\en{
    \item 
    Let $A, B$ be two $m \times n$ matrices. 
    Assume that 
    \begin{equation*}
        AX = BX 
    \end{equation*}
    for all $n$-tuples $X$. 
    Show that $A = B$. 
    \begin{Solution}
        \begin{proof}
            Let $C = A - B$. 
            Then $CX = O$ for all $X$. 
            Take $X = E^j$ to be the $j$-th unit vector for $j = 1, \dots, n$. 
            Then $CE^j = C^j$ is the $j$-th column of $C$. 
            By assumption, $CE^j = O$ for all $j$ so $C = O$. 
        \end{proof}
    \end{Solution}
    \item 
    \item 
    \item 
    \item 
    Let $P: V \to V$ be a linear map such that $P^2 = P$. 
    Define $Q = I - P$. 
    \ea{
        \item Show that $Q^2 = Q$. 
        \item Show that $\Im P = \Ker Q$ and $\Ker P = \Im Q$. 
    }
    A linear map $P$ such that $P^2 = P$ is called a \textbf{projection}. 
    It generalizes the notion of projection in the usual sense. 
    \begin{Solution}
        \ea{
            \item 
            \begin{proof}
                \begin{equation*}
                    Q^2 = (I - P)^2 = I^2 - 2IP + P^2 = I - 2P + P = I - P. 
                \end{equation*}
            \end{proof}
            \item 
            \begin{proof}
                Let $v \in \Im P$ so $v = Pw$ for some $w$. 
                Then $Qv = QPw = 0$ because $QP = (I - P)P = P - P^2 = P - P = O$. 
                Hence $\Im P \subset \Ker Q$. 
                Conversely, let $v \in \Ker Q$ so $Qv = 0$. 
                Then $(I - P)v = 0$ so $v - Pv = 0$, and $v = Pv$, so $v \in \Im P$, whence $\Ker Q \subset \Im P$. 
                Therefore we conclude that $\Im P = \Ker Q$. 

                Similarly, we can show that $\Ker P = \Im Q$. 
                Let $v \in \Ker P$ so $Pv = 0$. 
                Then $(I - Q)v = 0$ so $v - Qv = 0$, and $v = Qv$, so $v \in \Im Q$, whence $\Ker P \subset \Im Q$. 
                Conversely, let $v \in \Im Q$ so $v = Qw$ for some $w$. 
                Then $Pv = PQw = 0$ because $PQ = P(I - P) = P - P^2 = P - P = O$. 
                Hence $\Im Q \subset \Ker P$. 
                Therefore we conclude that $\Ker P = \Im Q$. 
            \end{proof}
        }
    \end{Solution}
   \item  
   Let $P: V \to V$ be a \textbf{projection}, that is the linear map such that $P^2 = P$. 
   \ea{
    \item 
    Show that $V = \Ker P + \Im P$. 
    \item 
    Show that the intersection of $\Im P$ and $\Ker P$ is $\{ O \}$. 
    In other words, if $v \in \Im P$ and $v \in \Ker P$, then $v = O$. 
   }
   \begin{Solution}
   \ea{
    \item 
    \begin{proof}
        Let $v \in V$. 
        Then $v = v - Pv + Pv$, and $v - Pv \in \Ker P$ because 
        \begin{equation*}
            P(v - Pv) = Pv - P^2v = Pv - Pv = O. 
        \end{equation*}
        Furthermore $Pv \in \Im P$, thus proving (a). 
    \end{proof}
    \item 
    \begin{proof}
        Let $v \in \Im P \cap \Ker P$. 
        Since $v \in \Im P$ there exists $w \in V$ such that $v = Pw$. 
        Since $v \in \Ker P$, we get 
        \begin{equation*}
            O = Pv = P^2w = Pw = v, 
        \end{equation*}
        so $v = O$, whence (b) is also proved. 
    \end{proof}
   }
   \end{Solution}

   Let $V$ be a vector space, and let $U, W$ be subspaces. 
   One says that $V$ is a \textbf{direct sum} of $U$ and $W$ if the following conditions are satisfied: 
   \begin{equation*}
        V = U + W \quad \text{and} \quad U \cap W = \{ O \}. 
   \end{equation*}
   In Exercise 6, we have proved that $V$ is the direct sum of $\Ker P$ and $\Im P$. 
   \item 
   \item 
   \item 
   Let $A, B$ be two matrices so that the product $AB$ is defined. 
   Prove that 
   \begin{equation*}
        \rank AB \leq \rank A \quad \text{and} \quad \rank AB \leq \rank B. 
   \end{equation*}
   [\textit{Hint}: Consider the linear maps $L_{AB}, L_A$, and $L_B$.] 
   \begin{solution}
       \begin{proof}
           Let $L_{AB}, L_A$ and $L_B$ be the linear maps such that
           \begin{equation*}
                L_{AB}(X) = ABX, \quad L_A(X) = AX \quad \text{and} \quad L_B(X) = BX. 
           \end{equation*}
           By composition of linear maps corresponding to multiplication of matrices, we have 
           \begin{equation*}
                L_{AB} = L_A \circ L_B. 
           \end{equation*}
           The rank of the matrix is the dimension of the image of its linear map, that is to say 
           \begin{equation*}
                \rank AB = \dim \Im L_{AB}, \quad \rank L_A = \dim \Im A \quad \text{and} \quad \rank B = \dim \Im B. 
           \end{equation*}

           We first show that $\rank AB \leq \rank A$. 
           It is clear that $\Im L_{AB} \subseteq \Im L_{A}$. 
           Then 
           \begin{equation*}
                \dim \Im L_{AB} \leq \dim \Im L_{A}, 
           \end{equation*}
           because the dimension of the subspace of a space is less than or equal to the dimension of that space. 
           Then we can immediately conclude that 
           \begin{equation*}
                \rank AB \leq \rank A. 
           \end{equation*}

           Next, we show that $\rank AB \leq \rank B$. 
           We have already shown before that the dimension of the image of a linear map is less than or equal to the dimension of its domain, say  
           \begin{equation*}
                \dim \Im L_{AB} = \dim \Im (L_A \circ L_B) \leq \dim \Im L_B, 
           \end{equation*}
           because we can regard $L_A$ as being restricted to the subspace $\Im L_B$ in the composition $L_{AB} = L_A \circ L_B$. 
           Then we sort the expressions and then conclude that 
           \begin{equation*}
                \rank AB \leq \rank B. 
           \end{equation*}
           Hence we complete the proof. 
       \end{proof} 
   \end{solution}
   In fact, define the \textbf{rank of a linear map} $L$ to be $\dim \Im L$. 
   If $L: V \to W$ and $F: W \to U$ are linear maps of finite dimensional spaces, show that 
   \begin{equation*}
        \rank F \circ L \leq \rank F \quad \text{and} \rank F \circ L \leq \rank L. 
   \end{equation*}
   \begin{Solution}
        \begin{proof}
            The dimension of a subspace is $\leq$ the dimension of the space. 
            Then 
            \begin{equation*}
                \Im F \circ L \leq \Im F \quad \text{so} \quad \dim \Im F \circ L \leq \dim \Im F 
            \end{equation*}
            so $\rank F \circ L \leq \rank F$. 
            This proves one of the formulas. 
            For the other, view $F$ as a linear map defined on $\Im L$. 
            Then 
            \begin{equation*}
                \rank F \circ L = \dim \Im F \circ L \leq \dim \Im L = \rank L. 
            \end{equation*}
        \end{proof}
   \end{Solution}
   \item 
}}